\scanpage{288}
\begin{center}- 102 -\end{center}\hfill IX

Or par le théorème d'orthogonalité appliqué a la fois pour $A_K$ et $A'_K$ on a
\[
V^\perp=W\quad,\quad(U/V)^\vee\simeq W'(-1)\quad,
\]
où $W'$ est la partie torique de $U'=T_\ell(A'(\overline K))$, et où le $(-1)$ désigne le twist de Tate. D'autre part, l'homomorphisme (9.2.4) se factorise par $I(\ell)\simeq\mathbb Z_\ell(1)$, de sorte que l'homomorphisme (9.2.4) peut s'interpréter comme un homomorphisme canonique
\[
(9.2.5)\qquad u:\mathbb Z_\ell(1)\longrightarrow W'\otimes W(-1)\quad,
\]
ou encore comme un élément
\[
(9.2.6)\qquad u\in W'\otimes W(-2)\quad.
\]
Se rappelant de la définition de $W$ et de $W'$, comme les groupes de points à valeurs dans $\overline K$ des deux faisceaux $\ell$-adiques constants tordus $T_\ell(A^o)^t$ et $T_\ell(A'^o)^t$, dont les fibres spéciales sont respectivement $T_\ell(T_o)$ et $T_\ell(T'_o)$, on trouve les isomorphismes canoniques
\[
(9.2.7)\quad\left\{\begin{array}{l}
W\simeq T_\ell(T_o)(\overline k))\simeq\check M\otimes\mathbb Z_\ell(1)\\[1.5ex]
W'\simeq T_\ell(T_o(\overline k))\simeq M'\otimes\mathbb Z_\ell(1)\quad,
\end{array}\right.
\]
$M$ et $M'$ désignant les groupes des caractères relativement à $\overline k$ de $T_o$ et de $T'_o$ respectivement :
\[
(9.2.8)\qquad T_o=D(M)\quad,\quad T'_o=D(M')\quad.
\]
Par suite, l'élément $u$ de (9.2.6), qui exprime l'opération du groupe d'inertie $I$ sur le module de Tate $U$ de $A_K$, s'interprète comme un élément

\margin{420}
\margin{420}
\margin{289}
